\newpage
\section{Praxisbeispiel: Der integrierte PP-Prozess bei Global Bike}

Dieses Kapitel überträgt die theoretischen Konzepte aus Kapitel 2 auf den integrierten Fertigungsprozess der Modellfirma Global Bike. Nach der Vorstellung des Szenarios wird der Prozess schrittweise durchlaufen -- von der Stammdatenvorbereitung über Grob- und Bedarfsplanung bis zu Fertigungsausführung und controllingseitigem Abschluss -- und abschließend kritisch gewürdigt.

\subsection{Vorstellung der Modellfirma und des Szenarios}
Global Bike ist das Modellunternehmen des Curriculums -- ein internationaler Fahrradhersteller mit den Gesellschaften Global Bike Inc. (USA) und Global Bike Germany GmbH.\footcites[Vgl.][Folie 8]{SAPUCC.2023}[][S. 2]{Boldau.2023} SAP S/4HANA bildet es hierarchisch ab: Der Mandant (Global Bike) ist die oberste Ebene, darunter der Buchungskreis als kleinste rechnungslegungspflichtige Einheit (US00 bzw. DE00); die Produktion erfolgt in Werken -- hier Dallas (DL00) -- mit Lagerorten für Roh-, Halb- und Fertigware.\footcite[Vgl.][Folie 7 ff.]{SAPUCC.2023}

Betrachtungsobjekt ist das Fertigerzeugnis Deluxe Touring Bike (DXTR, in Schwarz, Silber und Rot), das in Dallas in diskreter Fertigung montiert wird. Seine mehrstufige Produktstruktur umfasst einen farbspezifischen Rahmen, Räder und weitere Komponenten, die über Montagevorgänge zum Endprodukt zusammengeführt werden.\footcites[Vgl.][S. 4]{Boldau.2023}[][Folie 13 ff.]{SAPUCC.2023} Für die Grobplanung sind die Fahrräder zu einer Produktgruppe zusammengefasst.\footcite[Vgl.][S. 12]{Boldau.2023} Getragen wird der Prozess arbeitsteilig von mehreren Rollen -- u. a. Fertigungsleiter (Jun Lee), Werksleiter (Hiro Abe) und Controller (Jamie Shamblin) -- aus den Bereichen Produktionsplanung (PP) und Materialwirtschaft (MM).\footcite[Vgl.][S. 2]{Boldau.2023}

Das Szenario operationalisiert damit die Konzepte aus Kapitel 2: die Ebenen Mandant/Buchungskreis/Werk/Lagerort konkretisieren die Organisationsstruktur, die Produktgruppe die aggregierte SOP-Ebene (Kap. 2.2), der mehrstufige Aufbau des Fahrrads die Stückliste (Kap. 2.3) und seine Lagerfertigung die Strategie \enquote{Planung mit Endmontage} (Strategiegruppe 40, Kap. 2.6). Der folgende Abschnitt beginnt daher mit der Vorbereitung der Stammdaten.

\subsection{Vorbereitung der Stammdaten}
Bevor geplant werden kann, sind die Stammdaten des Deluxe Touring Bike um planungsrelevante Angaben zu erweitern. In der Rolle Fertigungsleiter wird der werksspezifische Materialstammsatz (DL00) angepasst: In der Dispositionssicht wird die Strategiegruppe 40 (Vorplanung mit Endmontage) gesetzt und damit die Make-to-Stock-Strategie aus Kap. 2.6 verankert; in der Prognosesicht werden zwölf Initialisierungsperioden hinterlegt, in den Steuerungsdaten der Optimierungsgrad \enquote{F (Fein)} mit Parameteroptimierung gewählt und die Glättungsfaktoren gesetzt: Alpha 0,20 (Grundwert), Beta 0,10 (Trend), Gamma 0,30 (Saison) und Delta 0,30 (MAD, mittlere absolute Abweichung).\footcite[Vgl.][S. 6]{Boldau.2023} Diese Werte steuern die exponentielle Glättung der Absatzprognose (Kap. 2.2) und werden analog für die silberne und schwarze Variante gepflegt.

Anschließend wird der Arbeitsplan angepasst. Er ist über Arbeitsplangruppe und Plangruppenzähler definiert, referenziert das Material und enthält je Vorgang Vorgabewerte und Zeitelemente für die Terminierung.\footcite[Vgl.][S. 8 f.]{Boldau.2023} Kern der Anpassung ist die Zuordnung (Allokation) der Stücklistenkomponenten zu den Vorgängen: Rahmen und Sitz zu Vorgang 0020, Lenker zu 0030, Aluminiumrad und Kettenschaltung zu 0040, Kette (0050), Bremsanlage (0060), Pedale (0070) sowie Garantiedokument und Verpackung zu 0100. Damit ist festgelegt, welche Komponente in welchem Fertigungsschritt verbaut wird -- ein abhängiger Prozess, bei dem jeder Vorgang auf dem vorhergehenden aufsetzt.\footcite[Vgl.][S. 9 f.]{Boldau.2023}

Diese Vorbereitung bestätigt die in Kap. 2.3 hergeleitete Rolle der Stammdaten: Erst die Verknüpfung der Dispositions- und Prognosedaten des Materialstammsatzes mit Stückliste und Arbeitsplan schafft die konsistente Datenbasis, ohne die weder Bedarfs- noch Terminplanung rechnen kann.

\subsection{Absatz- und Produktionsgrobplanung (SOP)}
Im vierten Schritt legt der Fertigungsleiter für die Produktgruppe (PG-DXTR, DL00) einen zwölfmonatigen Absatz- und Produktionsgrobplan (SOP) an, der Plandaten konsolidiert und künftige Mengen prognostiziert. Grundlage ist der historische Verbrauch, in der Fallstudie vorgegeben als Vergangenheitswerte von Mai 2017 bis März 2021; über die automatische Modellauswahl erkennt das System Trend und Saison und wendet ein Saison-Trend-Modell an, womit die Prognosemodelle aus Kap. 2.2 praktisch umgesetzt werden.\footcite[Vgl.][S. 14 ff.]{Boldau.2023} Aus dem Absatzplan wird ein absatzsynchroner Produktionsplan abgeleitet und um eine Zielreichweite (fünf Perioden) ergänzt, sodass die Produktionsmengen den Absatz decken und den gewünschten Lagerbestand sichern.\footcite[Vgl.][S. 16 f.]{Boldau.2023}

Da der SOP langfristig und aggregiert plant, enthält er noch keine diskreten Materialbedarfe; diese entstehen erst mit der Übergabe an die Programmplanung, die die Produktgruppenplanung auf die einzelnen Materialien herunterbricht (Übergabestrategie \enquote{Produktionsplan Material(ien) als Anteil PG}).\footcite[Vgl.][S. 18]{Boldau.2023} Dabei entstehen Planprimärbedarfe für die drei Fahrräder, aufgeteilt nach hinterlegtem Anteil -- DXTR1 40\,\%, DXTR2 30\,\%, DXTR3 30\,\%.\footcite[Vgl.][S. 19]{Boldau.2023} Damit ist der Übergang von der Grobplanung (Kap. 2.2) zur materialgenauen Feinplanung vollzogen; der Planprimärbedarf ist Ausgangspunkt der Bedarfsplanung.

\subsection{Programmplanung und Bedarfsplanung (MPS/MRP)}
Vor dem Planungslauf prüft der Werksleiter in der Programmplanung, ob für die drei Fahrräder Planprimärbedarfe (geplante unabhängige Bedarfe) vorliegen.\footcite[Vgl.][S. 20 f.]{Boldau.2023} Anschließend startet der Fertigungsleiter den Lauf aus Leitteileplanung (MPS) und Materialbedarfsplanung (MRP): Die MPS erzeugt Planaufträge nach den Vorgaben aus SOP und Programmplanung, parallel entstehen über die Stücklistenauflösung Planaufträge für die Sekundärbedarfe. Gesteuert wird der Lauf über Parameter, maßgeblich den Verarbeitungsschlüssel NETCH (Net-Change im gesamten Horizont), der nur geänderte Materialien neu plant, ergänzt um Vorgaben zu Bestellanforderung, Dispositionsliste, Planungsmodus und Eckterminierung.\footcite[Vgl.][S. 22]{Boldau.2023}

Inhaltlich führt das System die in Kap. 2.4 beschriebene Nettobedarfsrechnung durch: Bestand und feste Zugänge werden dem Sicherheitsbestand und den Bedarfen gegenübergestellt; ergibt sich eine Unterdeckung (dispositiv verfügbare Menge kleiner als null), legt die MRP Beschaffungsvorschläge -- Planaufträge bzw. Bestellanforderungen -- in der durch das Losgrößenverfahren bestimmten Menge an.\footcite[Vgl.][S. 23]{Boldau.2023}

Das Ergebnis zeigt die dynamische Bedarfs-/Bestandsliste: Für das rote Fahrrad weist sie zunächst keinen Bestand und keine frei verfügbare Menge aus; die Einträge lassen sich zu Periodensummen aus Planprimärbedarfen, geplanten Zugängen und ATP-Mengen (Available-to-Promise, frei verfügbare Mengen) verdichten.\footcite[Vgl.][S. 25 f.]{Boldau.2023} Über den Bedarfsverursacher wird sichtbar, dass der erste Planauftrag den Sicherheitsbestand und den ersten Planprimärbedarf deckt. Damit ist die Kette von der Grobplanung bis zum konkreten Planauftrag (Kap. 2.4) geschlossen.\footcite[Vgl.][S. 27]{Boldau.2023}

\subsection{Fertigungsausführung}
Zur Ausführung wird der Planauftrag aus der Bedarfs-/Bestandsliste in einen Fertigungsauftrag umgewandelt. Dabei terminiert das System, prüft die Verfügbarkeit, reserviert die Komponenten laut Stückliste, gibt den Auftrag frei und berechnet die Plankosten.\footcite[Vgl.][S. 28 f.]{Boldau.2023} Damit die Fertigung starten kann, werden die zuvor leeren Komponentenbestände über einen Wareneingang ins Lager aufgefüllt -- in der Fallstudie vereinfacht ohne den vorgelagerten Beschaffungsprozess.\footcite[Vgl.][S. 31]{Boldau.2023}

Anschließend bucht der Lagerarbeiter den Warenausgang der Komponenten mit Bezug auf den Fertigungsauftrag über die Bewegungsart 261 (Verbrauch für Auftrag): Die reservierten Materialien werden aus den Lagerorten entnommen -- das Aluminiumrad aus dem Halbfabrikatelager (SF00), die übrigen aus dem Rohstofflager (RM00) --, dem Auftrag als Ist-Kosten zugeordnet und über einen Material-, Buchhaltungs- und Kostenrechnungsbeleg fortgeschrieben.\footcite[Vgl.][S. 34 f.]{Boldau.2023} Nach der Montage meldet der Fertigungsarbeiter die Fertigstellung zurück (Endrückmeldung mit Reservierungsausbuchung und Gutmenge); das System berechnet die Fertigungskosten, und der Auftragsstatus wechselt von \enquote{freigegeben} zu \enquote{rückgemeldet}.\footcite[Vgl.][S. 40 ff.]{Boldau.2023}

Den Abschluss bildet der Wareneingang des fertigen Fahrrads in das Fertigerzeugnislager (FG00), mit dem der Wert des hergestellten Materials in den Auftrag fortgeschrieben wird.\footcite[Vgl.][S. 44 f.]{Boldau.2023} Der Ablauf entspricht damit exakt der in Kap. 2.5 hergeleiteten Logik aus Auftragsumwandlung, bewegungsartengesteuerten Warenbewegungen, Rückmeldung und Belegerzeugung.

\subsection{Controllingseitiger Abschluss}
Nach der Ausführung prüft der Controller in der Fertigungskostenanalyse die dem Auftrag zugeordneten Kosten; die Übersicht stellt summierte Soll- und Ist-Kosten gegenüber und weist Abweichungen aus (die Gemeinkostenzuschläge erscheinen nur in den Soll-Kosten).\footcite[Vgl.][S. 46 f.]{Boldau.2023} Im letzten Schritt rechnet der Controller den Fertigungsauftrag ab (Istabrechnung im Kostenrechnungskreis NA00, Periode gleich laufender Monat): Die zunächst nur temporär erfassten Kosten werden einem Kostenobjekt zugewiesen und der Auftrag entlastet.\footcite[Vgl.][S. 48]{Boldau.2023} Der Lauf erfolgt zunächst als Testlauf mit dem Bericht \enquote{Ist/Plan/Abweichung}, bevor er als Echtlauf gebucht wird.\footcite[Vgl.][S. 49 ff.]{Boldau.2023}

Damit ist der in Kap. 2.5 beschriebene controllingseitige Abschluss vollzogen: Die Ist-Kosten werden abgerechnet und über den Soll-Ist- bzw. Plan-Ist-Vergleich transparent gemacht, sodass Abweichungen erkennbar werden und der Kreis von der Planung über die Ausführung bis zur kostenrechnerischen Bewertung geschlossen ist.

\subsection{Kritische Würdigung des Prozesses}
Der Prozess weist deutliche Stärken auf. Zentral ist die durchgängige Integration: Planungs-, Ausführungs- und Abrechnungsdaten liegen auf einer gemeinsamen Datenbasis, sodass jede Warenbewegung automatisch einen Material-, Finanz- und Kostenrechnungsbeleg erzeugt und die Datenkonsistenz sichert.\footcites[Vgl.][S. 212 ff.]{Kurbel.2021}[][Folie 56]{SAPUCC.2023} Die Bedarfsableitung ist zudem weitgehend automatisiert -- aus dem Planprimärbedarf entstehen über Stücklistenauflösung und Nettobedarfsrechnung ohne manuelle Zwischenschritte Planaufträge (Kap. 3.4) --, und der Ablauf bleibt nachvollziehbar, da sich jeder Bedarf bis zum Verursacher zurückverfolgen lässt.

Diesen Stärken stehen Grenzen gegenüber. Erstens beruht die Planung auf umfangreichen, korrekt zu pflegenden Stammdaten und Parametern -- Strategiegruppe, Glättungsfaktoren, MRP-Steuerungsparameter (Kap. 3.2 und 3.4) --, was hohen Pflegeaufwand und Parametrisierungskomplexität bedeutet. Zweitens folgt der Ablauf einer Sukzessivplanung, deren Teilpläne auf Annahmen -- etwa über den künftigen Primärbedarf -- beruhen, die durch ungeplante Ereignisse verletzt werden können.\footcite[Vgl.][S. 33]{Kurbel.2021} Drittens ist die Prognose modellabhängig und grundsätzlich fehlerbehaftet, sodass Über- oder Unterdeckungen möglich bleiben.\footcite[Vgl.][Folie 27]{SAPUCC.2023}

Bezogen auf die Problemstellung aus Kapitel 1 löst der integrierte Ansatz das Abstimmungsproblem aus Nachfrage, Kapazität und Materialverfügbarkeit weitgehend: Er überführt eine Absatzprognose konsistent in konkrete Fertigungs- und Beschaffungsvorschläge und reduziert Intransparenz und isolierte Einzelplanung. Die Güte dieser Lösung hängt jedoch maßgeblich von Datenpflege und Planungsannahmen ab; der Prozess verlagert das Abstimmungsproblem damit von der manuellen Koordination hin zu einer daten- und parametergetriebenen Steuerung.
